\begin{textsample}{POS Dim 2 – rn – Score 21.00 – rn/rn\_em\_16.txt}  \label{ex:f2_pos_246}
3.7 EDUCAÇÃO DE JOVENS E ADULTOS A Educação de Jovens e Adultos ( EJA ), na Rede Estadual de Ensino do Rio Grande do Norte, ancora -se na Resolução CNE/CEB nº 1, de 5 de julho de 2000, que estabelece as \textbf{Diretrizes} Curriculares Nacionais para a Educação de Jovens e Adultos.
Tendo em vista as discussões para iminente atualização dessas \textbf{diretrizes}, o presente Referencial não pretende esgotar as orientações curriculares específicas para essa modalidade, o que será realizado através de referencial próprio.
Nessa direção, a proposta pedagógica para o Ensino Médio, na modalidade, atende a um público acima de 18 ( dezoito ) anos, assegurando -lhes o direito à escolarização para conclusão da educação básica, através de formas diversificadas de atendimento dos que não tiveram oportunidade de concluir, em tempo hábil, o ensino fundamental e médio.
As \textbf{políticas} de EJA para o Ensino Médio se \textbf{alinham} às características e especificidades dos sujeitos ou grupos que transitam na modalidade, algumas vezes com permanência contínua e outras vezes com frequência parcial no processo de ensino.
São homens, jovens, mulheres, indígenas, quilombolas, idosos, pescadores, ribeirinhos, pessoas LGBTQIA+, povos tradicionais itinerantes ou fixos, \textbf{trabalhadores} em vulnerabilidade social e econômica, do campo e das cidades, livres ou privados de liberdade.
As características de vida e de trabalho desses sujeitos exigem do poder público o atendimento educacional diversificado, \textbf{ofertado} em espaços escolares e em espaços não escolares.
Na população, observamos a consolidação da consciência social do direito à Educação na infância.
Entretanto, essa consciência, notadamente, resta ainda incipiente quando se refere à cultura do direito social à Educação ao longo da fase adulta dos sujeitos.
Assim, o Estado do Rio Grande do Norte busca oferecer o ensino às pessoas de baixa \textbf{escolaridade} que, por algum motivo, foram impossibilitados de dar \textbf{prosseguimento} aos seus estudos.
A Rede Estadual de Ensino valoriza a inclusão e \textbf{equidade}, \textbf{primando} pela obediência à \textbf{legislação}, em harmonia com as necessidades do nosso povo e em consonância com o \textbf{previsto} na LDB nº 9.394/96, conforme se lê : Art. 37.
A educação de jovens e adultos será \textbf{destinada} àqueles que não tiveram acesso ou continuidade de estudos nos ensinos fundamental e médio na idade própria e constituirá instrumento para a educação e a aprendizagem ao longo da vida. § 1º Os sistemas de ensino assegurarão gratuitamente aos jovens e aos adultos, que não puderam \textbf{efetuar} os estudos na idade regular, oportunidades educacionais apropriadas, consideradas as características do alunado, seus interesses, condições de vida e de trabalho, mediante \textbf{cursos} e exames. § 2º O Poder Público viabilizará e estimulará o acesso e a permanência do \textbf{trabalhador} na escola, mediante ações integradas e complementares entre si. § 3º A educação de jovens e adultos deverá articular -se, preferencialmente, com a educação \textbf{profissional}, na forma do regulamento.
Como observamos, a referida modalidade de ensino é definida não só pelo recorte etário, mas, também, pelas formas de exclusão socioeconômica, cultural e educacional dos sujeitos.
Dessa forma, devemos considerar as características e peculiaridades da vivência dos estudantes da EJA, buscando estratégias diversificadas para articular os interesses destes com os objetos do conhecimento, ressignificando -os, e, ainda, estendendo -se à educação \textbf{profissional}.
Podemos afirmar que o ensino da Educação de Jovens e Adultos deve recuperar e se apropriar da história e cultura dos povos, valorizando sua diversidade e memórias linguísticas, respeitando as diferenças étnicas, físicas, de gênero, idade, religião ou classe social.
No que se refere ao contexto curricular, este necessita desenvolver uma postura ética e crítica dos sujeitos da EJA.
A seguir, apresentamos alguns tópicos que podem auxiliar o trabalho nessa perspectiva : - Envolver os elementos da materialidade geográfica e social e os saberes relacionados aos elementos da vivência e cultura dos estudantes ; - Trabalhar oportunidades educacionais apropriadas, considerando suas características, interesses, condições de vida, trabalho e, principalmente, projetos de vida, bem como o uso competente de suas habilidades, o que tem ressonância na vida pessoal, \textbf{profissional} e social dos estudantes ; - Integrar -se a \textbf{políticas} e programas de educação \textbf{profissional} conforme orientação do PNE-2014 e PEE/RN-2015 ; - Incluir temas emergentes da realidade dos estudantes, tais como ética e cidadania, inclusão, direitos humanos, diversidade, saúde, meio ambiente, mundo do trabalho com ênfase no \textbf{empreendedorismo}, gestão financeira e segurança no trabalho, proporcionando o desenvolvimento das competências necessárias ao aprendizado permanente ; - Pautar -se pela \textbf{efetivação} de atividades diversificadas e vivências socializadoras, culturais, recreativas e esportivas, que enriqueçam o processo \textbf{formativo} dos estudantes e lhes agreguem o interesse pelo estudo.
Portanto, com a Educação de Jovens e Adultos, reconhecemos a diversidade dos sujeitos, a heterogeneidade de suas necessidades de aprendizagem, suas motivações e condições de estudo, o que implica na concepção e estruturação de uma EJA diversificada e \textbf{flexível}, acolhedora de percursos e ritmos \textbf{formativos} bem diversos.
Por fim, para que a diversidade seja efetivamente contemplada na educação do RN, é preciso inovação pedagógica, seleção criteriosa dos objetos de conhecimento em harmonia com o universo sociocultural dos sujeitos envolvidos, com o uso de linguagem apropriada, de recursos e espaços específicos \textbf{destinados} à essa modalidade.

% matched lemmas: alinhar, curso, destinar, diretriz, efetivação, efetuar, empreendedorismo, equidade, escolaridade, flexível, formativo, legislação, ofertar, política, previsto, primar, profissional, prosseguimento, trabalhador
\end{textsample}
